\chapter{Visão geral e responsabilidades}

Este manual destina-se à pessoa responsável por instalar, atualizar, proteger e recuperar o serviço em Ubuntu com systemd. O administrador precisa de acesso ao diretório do projeto, ao usuário do serviço e, para operações do sistema, a \code{sudo}.

\begin{noticebox}[Limite de suporte]
A aplicação foi planejada para rede interna, mas deve ser operada com controles equivalentes aos de produção. Não publique a porta diretamente na Internet. Use proxy reverso com HTTPS, firewall e política institucional de acesso.
\end{noticebox}

\section{Arquitetura}

O navegador carrega uma interface estática e chama a API HTTP do processo Node.js. A API autentica tokens e serializa mutações antes de gravar JSON por substituição atômica. Não há banco relacional nem dependência de runtime além do Node.js.

\begin{figure}[H]
\centering
\begin{tikzpicture}[
  archbox/.style={
    draw=Line,
    line width=0.6pt,
    rounded corners=2pt,
    fill=Paper,
    text width=31mm,
    minimum height=19mm,
    inner sep=2mm,
    align=center
  },
  serverbox/.style={archbox,draw=UtfprYellow!75!Ink,fill=UtfprYellow!14,line width=0.9pt},
  servicebox/.style={archbox,text width=38mm,minimum height=16mm},
  flow/.style={Latex-Latex,line width=0.8pt,color=Slate},
  supervision/.style={-{Latex[length=2mm]},line width=0.8pt,color=Slate},
  flowlabel/.style={fill=white,inner sep=1.5pt,font=\scriptsize,align=center,text=Slate}
]
\node[archbox] (browser) {%
  {\scriptsize\bfseries\color{Slate} CLIENTE}\\[1.2mm]
  {\bfseries Navegador}\\[-0.2mm]
  {\small HTML, CSS e JavaScript}%
};
\node[serverbox,right=27mm of browser] (node) {%
  {\scriptsize\bfseries\color{Slate} APLICAÇÃO}\\[1.2mm]
  {\bfseries Servidor Node.js}\\[-0.2mm]
  {\small HTTP + API autenticada}%
};
\node[archbox,right=27mm of node] (data) {%
  {\scriptsize\bfseries\color{Slate} PERSISTÊNCIA}\\[1.2mm]
  {\bfseries Diretório de dados}\\[-0.2mm]
  {\small JSON protegido}%
};

\draw[flow] (browser.east) --
  node[flowlabel] {requisições\\HTTPS/HTTP}
  (node.west);
\draw[flow] (node.east) --
  node[flowlabel] {leitura e escrita\\serializadas}
  (data.west);

\node[servicebox,below=14mm of node] (systemd) {%
  {\scriptsize\bfseries\color{Slate} SUPERVISÃO}\\[1.2mm]
  {\bfseries systemd}\\[-0.2mm]
  {\small ambiente e reinício}%
};
\draw[supervision] (systemd.north) --
  node[flowlabel,right=1mm] {inicia e\\monitora}
  (node.south);
\end{tikzpicture}
\caption{Limites da implantação e fluxos entre cliente, aplicação, persistência e supervisão. Somente a interface estática é publicada.}
\label{fig:admin-architecture}
\end{figure}

Somente \code{index.html}, folhas de estilo, JavaScript do navegador, logotipo e sprite de ícones pertencem à lista pública. Arquivos em \code{data/}, \code{.git/}, scripts, configuração e código do servidor devem sempre retornar negação ou não encontrado por HTTP, inclusive em caminhos codificados.

\chapter{Instalação em Ubuntu}

\section{Pré-requisitos}

Use um usuário dedicado ou técnico que continuará proprietário dos dados e executará o serviço. São necessários Ubuntu suportado, acesso administrativo, Git para atualização e Node.js 20 ou superior. Reserve um diretório estável, por exemplo \code{/opt/mapa-horarios}, e um diretório de dados incluído na rotina de backup.

\begin{enumerate}
  \item Obtenha o projeto pelo canal institucional autorizado.
  \item Entre no diretório e revise \code{setup.sh} antes da primeira execução.
  \item Execute o script como o mesmo usuário definido no serviço:
\end{enumerate}

\begin{lstlisting}
cd /opt/mapa-horarios
./setup.sh
\end{lstlisting}

O script verifica Node.js/npm, executa \code{npm ci --omit=dev}, protege o diretório de dados, cria \path{/etc/mapa-horarios.env} se ausente e instala \path{/etc/systemd/system/mapa-horarios.service}. O segredo gerado é preservado nas execuções posteriores.

\section{Primeiro usuário}

Depois da instalação, crie pelo menos um usuário e reinicie o serviço se necessário:

\begin{lstlisting}
./scripts/add-user.sh editor 'senha-longa-e-exclusiva'
sudo systemctl restart mapa-horarios
\end{lstlisting}

Teste localmente \url{http://127.0.0.1:3000}. O padrão escuta somente no endereço de loopback.

\section{Implantação e systemd}

\begin{figure}[H]
\centering
\begin{tikzpicture}[
  node distance=9mm,
  processbox/.style={draw=Line,rounded corners=2pt,fill=Paper,minimum width=65mm,minimum height=10mm,align=center},
  flow/.style={-{Latex[length=2mm]},thick,color=Slate}
]
\node[processbox] (env) {\code{/etc/mapa-horarios.env}\\segredo, host e dados};
\node[processbox,below=of env] (unit) {\code{mapa-horarios.service}\\usuário e diretório de trabalho};
\node[processbox,below=of unit] (process) {\code{node server.js}\\\code{127.0.0.1:3000}};
\node[processbox,below=of process] (proxy) {Proxy HTTPS/rede interna\\controle de acesso e firewall};
\draw[flow] (env) -- (unit);
\draw[flow] (unit) -- (process);
\draw[flow] (proxy) -- (process);
\end{tikzpicture}
\caption{Relação entre ambiente protegido, unidade systemd, processo e acesso de rede.}
\label{fig:admin-deployment}
\end{figure}

\chapter{Configuração segura}

\section{Variáveis de ambiente}

Mantenha as variáveis em \code{/etc/mapa-horarios.env}, com leitura restrita ao administrador e ao usuário do serviço.

\begin{longtable}{@{}p{0.27\textwidth}p{0.18\textwidth}p{0.49\textwidth}@{}}
\toprule
\textbf{Variável} & \textbf{Padrão} & \textbf{Finalidade} \\
\midrule
\endhead
\code{JWT\_SECRET} & nenhum em produção & Segredo de no mínimo 32 caracteres. A inicialização falha se estiver ausente ou fraco com \code{NODE\_ENV=production}. \\
\code{AUTH\_TOKEN\_TTL\_MS} & \code{3600000} & Duração da sessão em milissegundos. \\
\code{HOST} & \code{127.0.0.1} & Interface de rede em que o processo escuta. \\
\code{PORT} & \code{3000} & Porta HTTP do processo. \\
\code{MAPA\_DATA\_DIR} & \code{./data} & Diretório de usuários e configurações. \\
\bottomrule
\end{longtable}

Exemplo mínimo:

\begin{lstlisting}
JWT_SECRET=valor-aleatorio-com-pelo-menos-32-caracteres
HOST=127.0.0.1
PORT=3000
MAPA_DATA_DIR=/var/lib/mapa-horarios
AUTH_TOKEN_TTL_MS=3600000
\end{lstlisting}

Depois de editar o ambiente:

\begin{lstlisting}
sudo systemctl restart mapa-horarios
sudo systemctl --no-pager --full status mapa-horarios
\end{lstlisting}

\section{Autenticação}

O login confere um hash \code{scrypt} e emite um token assinado, armazenado no \code{sessionStorage} do navegador. Toda operação da biblioteca, inclusive leitura, exige \code{Authorization: Bearer}. Tokens expirados, adulterados ou malformados são recusados; respostas de login não revelam se o usuário existe.

\begin{figure}[H]
\centering
\begin{tikzpicture}[
  node distance=10mm and 13mm,
  box/.style={draw=Line,rounded corners=2pt,fill=Paper,minimum width=34mm,minimum height=10mm,align=center},
  flow/.style={-{Latex[length=2mm]},thick,color=Slate}
]
\node[box] (login) {Usuário + senha};
\node[box,right=of login] (verify) {Verificação scrypt\\e limite de tentativas};
\node[box,right=of verify] (token) {Token assinado\\com expiração};
\node[box,below=of token] (api) {API protegida};
\draw[flow] (login) -- (verify);
\draw[flow] (verify) -- (token);
\draw[flow] (token) -- node[right,font=\scriptsize] {Bearer} (api);
\end{tikzpicture}
\caption{Fluxo de autenticação. Credenciais e tokens nunca devem aparecer em logs ou documentação.}
\label{fig:admin-auth}
\end{figure}

\chapter{Usuários}

\section{Criar e remover}

Os scripts usam \code{data/users.json} por padrão ou o arquivo indicado com \code{--file}. Execute-os como o usuário do serviço para preservar propriedade e permissões.

\begin{lstlisting}
./scripts/add-user.sh <username> '<senha>'
./scripts/add-user.sh <username> '<senha>' --id user-identificador
./scripts/remove-user.sh <username-ou-id>
\end{lstlisting}

Não passe senhas em histórico compartilhado ou gravação de terminal. Após criar ou remover, confirme que o arquivo continua com modo \code{640} e diretórios com \code{750}.

\section{Redefinir senha}

Não existe comando de alteração direta. Para evitar ficar sem acesso, mantenha uma segunda conta administrativa válida, remova a conta antiga e adicione-a novamente com nova senha:

\begin{lstlisting}
./scripts/add-user.sh administrador-temporario '<senha-temporaria>'
./scripts/remove-user.sh editor
./scripts/add-user.sh editor '<nova-senha-forte>'
./scripts/remove-user.sh administrador-temporario
\end{lstlisting}

Faça backup de \code{users.json} antes da operação e teste a nova conta em uma sessão anônima.

\chapter{Dados, revisões e concorrência}

O diretório configurado contém \code{users.json} e \code{configurations.json}. Ambos são arrays JSON e dados de runtime; não devem ser adicionados ao Git. Os arquivos \code{*.example.json} vazios servem somente para instalações novas.

Cada configuração contém identificador, nome, \code{schemaVersion}, \code{revision}, datas, autores, estado da interface e contadores. Registros antigos são normalizados em memória e persistem os novos metadados no próximo salvamento válido.

\section{Revisão otimista e escrita atômica}

\begin{figure}[H]
\centering
\begin{tikzpicture}[
  node distance=9mm,
  processbox/.style={draw=Line,rounded corners=2pt,fill=Paper,text width=48mm,minimum height=9mm,align=center},
  flow/.style={-{Latex[length=2mm]},thick,color=Slate}
]
\node[processbox] (request) {Cliente envia \code{If-Match: "revision"}};
\node[processbox,below=of request] (compare) {Servidor compara com a revisão atual};
\node[processbox,below left=9mm and 2mm of compare] (conflict) {Divergente\\\code{409 Conflict}};
\node[processbox,below right=9mm and 2mm of compare] (queue) {Igual\\entra na fila};
\node[processbox,below=of queue] (atomic) {Grava temporário, substitui original\\e incrementa revisão};
\draw[flow] (request) -- (compare);
\draw[flow] (compare) -- (conflict);
\draw[flow] (compare) -- (queue);
\draw[flow] (queue) -- (atomic);
\end{tikzpicture}
\caption{Controle de concorrência. Sem o cabeçalho de revisão, uma atualização ou exclusão retorna o estado HTTP 428.}
\label{fig:admin-revision}
\end{figure}

Uma revisão divergente nunca é sobrescrita automaticamente. A interface permite recarregar a versão do servidor ou criar uma cópia. A fila serializa mutações no mesmo processo e a troca atômica evita JSON parcial. Esse modelo não é indicado para múltiplas instâncias compartilhando o mesmo diretório.

\chapter{Backup e restauração}

\section{Política mínima}

Faça backup dos dois arquivos, mantenha cópias fora do servidor e teste restaurações periodicamente. Proteja backups como credenciais, pois \code{users.json} contém hashes e as configurações podem conter informações acadêmicas.

\begin{figure}[H]
\centering
\begin{tikzpicture}[
  node distance=7mm,
  processbox/.style={draw=Line,rounded corners=2pt,fill=Paper,minimum width=75mm,minimum height=9mm,align=center},
  flow/.style={-{Latex[length=2mm]},thick,color=Slate}
]
\node[processbox] (stop) {Parar serviço};
\node[processbox,below=of stop] (copy) {Copiar os dois JSON para destino protegido};
\node[processbox,below=of copy] (validate) {Validar integridade e registrar data};
\node[processbox,below=of validate] (start) {Iniciar serviço e testar login/biblioteca};
\draw[flow] (stop) -- (copy);
\draw[flow] (copy) -- (validate);
\draw[flow] (validate) -- (start);
\end{tikzpicture}
\caption{Fluxo conservador para backup consistente ou restauração.}
\label{fig:admin-backup}
\end{figure}

\section{Criar backup consistente}

Substitua os caminhos pelos valores reais e use um diretório novo com data explícita:

\begin{lstlisting}
sudo systemctl stop mapa-horarios
sudo install -d -m 700 /srv/backups/mapa/2026-08-01
sudo cp /var/lib/mapa-horarios/configurations.json /srv/backups/mapa/2026-08-01/
sudo cp /var/lib/mapa-horarios/users.json /srv/backups/mapa/2026-08-01/
node -e "JSON.parse(require('fs').readFileSync(process.argv[1]))" \
  /srv/backups/mapa/2026-08-01/configurations.json
node -e "JSON.parse(require('fs').readFileSync(process.argv[1]))" \
  /srv/backups/mapa/2026-08-01/users.json
sudo systemctl start mapa-horarios
\end{lstlisting}

\section{Restaurar}

\begin{dangerbox}
A restauração substitui o estado atual. Antes dela, faça uma cópia de emergência dos arquivos existentes, mesmo quando estiverem possivelmente corrompidos.
\end{dangerbox}

\begin{lstlisting}
sudo systemctl stop mapa-horarios
sudo cp /var/lib/mapa-horarios/configurations.json /tmp/configurations.pre-restore.json
sudo cp /var/lib/mapa-horarios/users.json /tmp/users.pre-restore.json
sudo cp /srv/backups/mapa/2026-08-01/configurations.json /var/lib/mapa-horarios/
sudo cp /srv/backups/mapa/2026-08-01/users.json /var/lib/mapa-horarios/
sudo chown mapa:mapa /var/lib/mapa-horarios/*.json
sudo chmod 640 /var/lib/mapa-horarios/*.json
sudo systemctl start mapa-horarios
\end{lstlisting}

Troque \code{mapa:mapa} pelo usuário e grupo reais da unidade. Verifique status, logs, login, quantidade de semestres e abertura de um registro.

\chapter{Atualização, rotação e recuperação}

\section{Atualização segura}

\begin{enumerate}
  \item Anuncie a janela e confirme que não há edição em andamento.
  \item Registre o commit atual e faça backup dos dados e do arquivo de ambiente.
  \item Pare o serviço.
  \item Atualize o código pelo fluxo Git aprovado, sem apagar o diretório de dados.
  \item Execute \code{npm ci --omit=dev} e as verificações disponíveis em uma cópia limpa.
  \item Inicie o serviço, acompanhe logs e faça um teste funcional.
\end{enumerate}

\begin{lstlisting}
git rev-parse HEAD
sudo systemctl stop mapa-horarios
git pull --ff-only
npm ci --omit=dev
sudo systemctl start mapa-horarios
sudo journalctl -u mapa-horarios -n 100 --no-pager
\end{lstlisting}

Se o teste falhar, pare o processo, retorne ao artefato/commit anteriormente aprovado pelo procedimento institucional, restaure o backup somente se o formato dos dados tiver sido alterado e documente a ocorrência. Não use comandos destrutivos sobre a raiz do projeto.

\section{Rotacionar o segredo JWT}

Gere um novo valor fora do repositório, substitua \code{JWT\_SECRET} no arquivo protegido e reinicie. Todas as sessões existentes serão invalidadas.

\begin{lstlisting}
node -e "process.stdout.write(require('crypto').randomBytes(32).toString('hex'))"
sudoedit /etc/mapa-horarios.env
sudo systemctl restart mapa-horarios
\end{lstlisting}

Faça a rotação após suspeita de exposição, mudança de responsáveis ou política periódica. Não registre o valor em chamados, logs ou histórico do shell.

\section{Arquivos JSON inválidos}

Se o processo relatar JSON inválido:

\begin{enumerate}
  \item pare o serviço e preserve uma cópia do arquivo problemático;
  \item valide sintaxe sem imprimir o conteúdo em um canal compartilhado;
  \item restaure o backup íntegro mais recente;
  \item aplique propriedade e modo corretos;
  \item inicie e valide a biblioteca.
\end{enumerate}

\chapter{Operação, logs e diagnóstico}

\section{Comandos do serviço}

\begin{lstlisting}
sudo systemctl status mapa-horarios
sudo systemctl restart mapa-horarios
sudo systemctl stop mapa-horarios
sudo systemctl start mapa-horarios
sudo journalctl -u mapa-horarios -n 100 --no-pager
sudo journalctl -u mapa-horarios -f
\end{lstlisting}

O projeto também fornece \code{./stop.sh}. Prefira os comandos systemd durante diagnóstico para manter uma única fonte de estado.

\section{Problemas frequentes}

\begin{longtable}{@{}p{0.31\textwidth}p{0.63\textwidth}@{}}
\toprule
\textbf{Sintoma} & \textbf{Verificação} \\
\midrule
\endhead
Serviço não inicia em produção & Confira \code{JWT\_SECRET}, tamanho mínimo, permissões do arquivo de ambiente e logs. \\
Porta já utilizada & Use \code{ss -ltnp}, identifique o processo autorizado ou altere \code{PORT}; atualize também o proxy. \\
Login sempre falha & Valide \code{users.json}, propriedade, modo e relógio do servidor. Crie uma conta de teste temporária. \\
API retorna 401/403 & Sessão ausente, expirada ou token inválido; faça novo login. \\
API retorna 409 & Edição concorrente normal; preserve como cópia ou recarregue. \\
API retorna 428 & Cliente de API não enviou \code{If-Match} em PUT/DELETE. \\
Semestres desapareceram & Confira \code{MAPA\_DATA\_DIR}, usuário do serviço e montagem do volume antes de restaurar. \\
Permissão negada ao salvar & Confirme propriedade do diretório, modo \code{750} e arquivos \code{640}. \\
\bottomrule
\end{longtable}

\chapter{Verificação e checklists}

Para desenvolvimento e homologação:

\begin{lstlisting}
npm ci
npm run check
npx playwright install --with-deps chromium
npm run test:e2e
npm run docs:check
\end{lstlisting}

\code{npm run check} executa lint, verificação de ícones, busca por dados sensíveis e testes de servidor. Playwright cobre o navegador. A documentação verifica metadados, figuras e compilação dos dois PDFs.

\section{Após instalar ou atualizar}

\begin{itemize}
  \item[$\square$] Unidade ativa e sem reinícios repetidos.
  \item[$\square$] Processo ouvindo somente no host e porta planejados.
  \item[$\square$] HTTPS e firewall validados quando houver acesso pela rede.
  \item[$\square$] Login válido e inválido testados sem revelar detalhes.
  \item[$\square$] Biblioteca exige sessão e abre um semestre.
  \item[$\square$] Salvamento, exportação e pré-visualização testados.
  \item[$\square$] URLs de \code{data/}, \code{server.js} e \code{.git/} não expõem conteúdo.
  \item[$\square$] Backup recente, íntegro e recuperável registrado.
\end{itemize}

\section{Após restaurar}

\begin{itemize}
  \item[$\square$] Ambos os arquivos são arrays JSON válidos.
  \item[$\square$] Proprietário, grupo e modos conferidos.
  \item[$\square$] Serviço iniciou sem erros de leitura.
  \item[$\square$] Conta administrativa autenticou.
  \item[$\square$] Quantidade e conteúdo amostral de semestres conferidos.
  \item[$\square$] Cópia anterior à restauração mantida até a validação final.
\end{itemize}

\section{Resposta a incidente}

\begin{itemize}
  \item[$\square$] Preservar logs e evidências sem divulgar segredos.
  \item[$\square$] Restringir acesso ou parar o serviço se houver exposição ativa.
  \item[$\square$] Rotacionar JWT e redefinir credenciais afetadas.
  \item[$\square$] Verificar integridade de código, ambiente, usuários, dados e backups.
  \item[$\square$] Restaurar somente de fonte confiável e testar.
  \item[$\square$] Registrar causa, período, impacto e ação preventiva.
\end{itemize}
